Again put \(\delta=\delta_n\), \(h=h_n\), \(b_x=p_xq_x(\delta)\), and \(m_j=|I_j|\).  Consider the simultaneous events
\[
 \begin{split}
 E_0&=\{|\widehat\nu_n-\nu_\delta(P)|\leq b_0\},\\
 E_{1x}&=\{|\widehat b_x-b_x|\leq b_1\},\\
 E_{2x}&=\Omega_x^c\ \cup\
 \left\{\left|\sum_{i\in I_2}w_{ix}Y_i-\mu_x^P(\delta)\right|\leq B_x\right\}.
 \end{split} \tag{17}
\]
We first derive the concentration inequality used below.  If a centered random variable \(W\) lies in an interval of length \(d\), its cumulant generating function \(K(\lambda)=\log\mathbb E e^{\lambda W}\) satisfies \(K(0)=K'(0)=0\) and
\[
 K''(\lambda)=\operatorname{Var}_{\lambda}(W)\leq d^2/4,
\]
because variance under any exponentially tilted law is at most the second moment about the midpoint of that interval.  Two integrations give \(K(\lambda)\leq\lambda^2d^2/8\).  Independence, exponential Markov, and optimization over \(\lambda>0\) therefore yield, for independent \(Z_i\in[0,1]\) and deterministic \(a_i\),
\[
 P\left(\left|\sum_i a_i\{Z_i-\mathbb EZ_i\}\right|>
 t\Bigl(\sum_i a_i^2\Bigr)^{1/2}\right)\leq2e^{-2t^2}. \tag{18}
\]
This derivation also applies conditionally when the weights and conditional means are fixed.

Apply (18) with equal weights to the first two sample blocks.  The definition \(t_\alpha=\sqrt{\log(12J/\alpha)/2}\) gives
\[
 P(E_0^c)\leq2e^{-2t_\alpha^2}=\frac{\alpha}{6J},
 \qquad
 P(E_{1x}^c)\leq2e^{-2t_\alpha^2}=\frac{\alpha}{6J}. \tag{19}
\]
On \(\Omega_x\), condition on \(\mathscr D_2=\sigma\{(X_i,A_i):i\in I_2\}\).  Assumption \texttt{ass:iid-sampling} gives conditional independence.  Applying (18) to the weighted outcome residuals yields
\[
 P\left(\left.\left|\sum_iw_{ix}\{Y_i-\mu_x^P(A_i)\}\right|>
 t_\alpha\Bigl(\sum_iw_{ix}^2\Bigr)^{1/2}\,\right|\mathscr D_2\right)
 \leq\frac{\alpha}{6J}. \tag{20}
\]
By Lemma \texttt{lem:total-gram}, the weights reproduce every polynomial through degree \(\ell\).  Assumption \texttt{ass:holder-regression} consequently bounds the conditional bias by \(Lh^\beta\sum_i|w_{ix}|\).  Thus (20) proves \(P(E_{2x}^c)\leq\alpha/(6J)\).  A union bound and fixed finiteness of \(J\) give
\[
 P\left(E_0\cap\bigcap_{x\in\mathcal X}(E_{1x}\cap E_{2x})\right)
 \geq1-\frac{(2J+1)\alpha}{6J}\geq1-\alpha. \tag{21}
\]

On the event in (21), projection onto \([0,1]\) cannot enlarge the local regression error.  On \(\Omega_x\), \(b_x\leq\widehat b_x+b_1\) and hence
\[
 |\widehat b_x\widetilde\mu_x-b_x\mu_x^P(\delta)|
 \leq b_1+(\widehat b_x+b_1)B_x=D_x. \tag{22}
\]
On \(\Omega_x^c\), the two products lie in \([0,\widehat b_x]\) and \([0,b_x]\), respectively, so
\[
 |\widehat b_x\widetilde\mu_x-b_x\mu_x^P(\delta)|
 \leq\max\{\widehat b_x,b_x\}\leq\widehat b_x+b_1=D_x. \tag{23}
\]
Combining (22)--(23) with \(E_0\), then using contraction of the final projection about \(\theta_\delta(P)\in[0,1]\), gives
\[
 |\widehat\theta_n-\theta_\delta(P)|\leq b_0+\sum_xD_x. \tag{24}
\]
By Definition \texttt{def:honest-interval}, (24) is precisely the event \(\theta_\delta(P)\in\mathrm{CI}_n\).  Equation (21) proves the asserted uniform finite-sample coverage.

It remains to integrate the radius, including singular-Gram realizations.  The binomial calculation (4)--(5) in the proof of Theorem \texttt{thm:minimax-risk}, Lemma \texttt{lem:total-gram}, and Lemma \texttt{lem:bandwidth-phases} imply
\[
 \begin{aligned}
 \mathbb E_P[B_x\mathbf 1\{\Omega_x\}]
 &\leq Lh^\beta\mathbb E_P\!\left[\sum_i|w_{ix}|\mathbf 1\{\Omega_x\}\right]
 +t_\alpha\mathbb E_P\!\left[\Bigl(\sum_iw_{ix}^2\Bigr)^{1/2}\mathbf 1\{\Omega_x\}\right]\\
 &\leq Ch^\beta+C\mathbb E_P[N_x^{-1/2}\mathbf 1\{N_x>0\}]
 \leq Ch^\beta. 
 \end{aligned} \tag{25}
\]
The split makes \(\widehat b_x\) independent of \((B_x,\Omega_x)\), and \(\mathbb E_P\widehat b_x=b_x\).  The two branches in Definition \texttt{def:honest-interval} therefore give
\[
 \begin{aligned}
 \mathbb E_PD_x
 &\leq(b_x+b_1)\mathbb E_P[B_x\mathbf 1\{\Omega_x\}]+b_1
 +(b_x+b_1)P(\Omega_x^c)\\
 &\leq C\{n^{-1/2}+b_xh^\beta\}
 \leq C\{n^{-1/2}+\delta^{\kappa+1}h^\beta\}. 
 \end{aligned} \tag{26}
\]
For the second line, Proposition \texttt{prop:pushforward-setup} gives \(b_x\leq C\delta^{\kappa+1}\), while Lemmas \texttt{lem:total-gram} and \texttt{lem:bandwidth-phases} give
\[
 P(\Omega_x^c)\leq Ce^{-cnh(\delta+h)^\kappa}
 \leq Ce^{-ch^{-2\beta}}\leq Ch^\beta. \tag{27}
\]
The terms containing \(b_1\) are \(O(n^{-1/2})\), since \(h\leq1\).  Intersecting with \([0,1]\) cannot increase length.  Thus (26), fixed \(J\), and \(b_0=O(n^{-1/2})\) yield
\[
 \sup_{P\in\mathcal M}\mathbb E_P[\operatorname{len}(\mathrm{CI}_n)]
 \leq2b_0+2\sum_x\mathbb E_PD_x\leq Cr_n. \tag{28}
\]

For the converse, use the global and localized same-class pairs explicitly constructed in Theorem \texttt{thm:minimax-risk}.  Decrease their fixed amplitudes, if necessary, so that their product-law total variation is at most a constant \(\eta<1-2\alpha\); equations (9), (12), and (14) of that proof permit this without changing either separation order.  For either pair, let \(Q_0,Q_1\) be the product laws and let their target values have separation \(\Delta\).  Uniform coverage of an arbitrary interval \(C_n\) gives
\[
 \begin{aligned}
 Q_0\{\theta_0,\theta_1\in C_n\}
 &\geq Q_0\{\theta_0\in C_n\}+Q_0\{\theta_1\in C_n\}-1\\
 &\geq(1-\alpha)+(1-\alpha-\eta)-1
 =1-2\alpha-\eta>0. 
 \end{aligned} \tag{29}
\]
On this event an interval has length at least \(\Delta\), whence
\[
 \sup_{P\in\mathcal M}\mathbb E_P[\operatorname{len}(C_n)]
 \geq\Delta(1-2\alpha-\eta). \tag{30}
\]
The global pair has \(\Delta\asymp n^{-1/2}\), and the localized pair has \(\Delta\geq c\delta^{\kappa+1}h^\beta\).  Maximizing the two applications of (30) and using \(\max(u,v)\geq(u+v)/2\) proves \(L_n^\star\geq cr_n\).  Together, (21), (28), and (30) prove uniform honesty, the achieved expected length, and its minimax lower bound.